%\input{tcilatex}


\documentclass{article}
%%%%%%%%%%%%%%%%%%%%%%%%%%%%%%%%%%%%%%%%%%%%%%%%%%%%%%%%%%%%%%%%%%%%%%%%%%%%%%%%%%%%%%%%%%%%%%%%%%%%%%%%%%%%%%%%%%%%%%%%%%%%%%
\usepackage{graphicx}
\usepackage{amsmath}

\begin{document}


\section{Flux-Surface and Bounce Geometry}

Three different coordinates system are considered here.

\begin{itemize}
\item  First, the toroidal coordinate system $\left( R,Z,\phi \right) $ with
local orthogonal vectors 
\begin{eqnarray}
\widehat{R} &=&\nabla R \\
\widehat{Z} &=&\nabla Z \\
\widehat{\phi } &=&R\nabla \phi 
\end{eqnarray}
This coordinate system conserves the largest generality in the magnetic
geometry

\item  Second, the poloidal (polar) coordinate system $\left( r,\theta ,\phi
\right) $, centered around the magnetic axis $\left( R_{0},Z_{0}\right) $
defined as the position of an extremum of $\psi $ (can be arbitrarily chosen
as $\psi =0$), and with local orthogonal vectors 
\begin{eqnarray}
\widehat{r} &=&\nabla r \\
\widehat{\theta } &=&r\nabla \theta  \\
\widehat{\phi } &=&R\nabla \phi 
\end{eqnarray}
and the relations 
\begin{eqnarray}
R &=&R_{0}+r\cos \theta  \\
Z &=&Z_{0}+r\sin \theta 
\end{eqnarray}
which can be inverted as 
\begin{eqnarray}
r &=&\sqrt{\left( R-R_{0}\right) ^{2}+\left( Z-Z_{0}\right) ^{2}} \\
\theta  &=&\arctan \left( \frac{Z-Z_{0}}{R-R_{0}}\right) +\pi \text{sign}%
\left( R-R_{0}\right)   \label{theta}
\end{eqnarray}

\item  Third, the flux coordinate system $\left( \psi ,s,\phi \right) $ is
the natural system for Fokker-Planck simulations, since the particles are
essentially confined to a given flux surface $\psi $. The distance $s$ is
the length along the poloidal magnetic field lines. We can choose its origin
as being at the position of minimum $B$-field amplitude within a
flux-surface. 
\begin{equation}
B\left( \psi ,s\equiv 0\right) =\min \left\{ B\left( \psi ,s\right) \right\}
=B_{\min }\left( \psi \right) 
\end{equation}
The direction is counter-clockwise and the limits $s_{\min }$ and $s_{\max }$
are set at the position of maximum magnetic field

\begin{equation}
B\left( \psi ,s\equiv s_{\min }\right) =B\left( \psi ,s\equiv s_{\max
}\right) =\max \left\{ B\left( \psi ,s\right) \right\} =B_{\max }\left( \psi
\right) 
\end{equation}
The local orthogonal vectors are 
\begin{eqnarray}
\widehat{\psi } &=&\frac{\nabla \psi }{\left| \nabla \psi \right| } \\
\widehat{s} &=&\frac{\nabla \phi \times \nabla \psi }{\left| \nabla \psi
\right| \left| \nabla \phi \right| }=\widehat{b}_{p} \\
\widehat{\phi } &=&R\nabla \phi 
\end{eqnarray}

\item  An alternative to the previous system may be easier to implement
numerically for the calculation of the bounce coefficients. This is the
system $\left( \psi ,\theta ,\phi \right) $. The $\theta $ grid is defined
in (\ref{theta}). We also define, for geometrical purposes, a flux-function $%
\rho \left( \psi \right) $ which coincides with the normalized radius on the
horizontal LFS midplane. Indeed, we have
\[
\rho \left( R>R_{0},Z=Z_{0}\right) =\frac{R-R_{0}}{R_{\max }-R_{0}}
\]
where $R_{\max }$ is the value of $R$ on the separatrix as it crosses the
midplane. The local orthogonal vectors are 
\begin{eqnarray}
\widehat{\psi } &=&\frac{\nabla \psi }{\left| \nabla \psi \right| } \\
\widehat{\theta } &=&\frac{\nabla \phi \times \nabla \psi }{\left| \nabla
\psi \right| \left| \nabla \phi \right| }=\widehat{s}=\widehat{b}_{p} \\
\widehat{\phi } &=&R\nabla \phi 
\end{eqnarray}
\end{itemize}

The 2-D outputs from HELENA\ are given on the $\left( \psi ,\theta \right) $
grid.

\section{Bounce Time}

\subsection{Normalized Expression}

The transit, or bounce time is defined as 
\begin{eqnarray}
\tau _{b}\left( \psi \right) &=&\int_{s_{\min }}^{s_{\max }}\frac{ds}{\left|
v_{s}\right| } \\
&=&\int_{s_{\min }}^{s_{\max }}\frac{ds}{\left| v_{\Vert }\right| }\frac{B}{%
\left| B_{P}\right| }
\end{eqnarray}

Using 
\begin{equation}
ds=Jd\theta 
\end{equation}
and 
\begin{equation}
v_{\Vert }=v\xi 
\end{equation}
we find 
\begin{equation}
\tau _{b}\left( \psi \right) =\frac{2\pi }{v\left| \xi _{0}\right| }%
\int_{\theta _{\min }}^{\theta _{\max }}\frac{d\theta }{2\pi }J\frac{\xi _{0}%
}{\xi }\frac{B}{\left| B_{P}\right| }
\end{equation}
where $\xi _{0}$ is the pitch angle at the position $\theta _{0}$ of minimum 
$B$-field. 
\begin{equation}
\theta _{0}\equiv \theta \left( B=B_{\min }\left( \psi \right) \right) 
\end{equation}
$\ $

The bounce time can be normalized as such: 
\begin{equation}
\tau _{b}\left( \psi ,\xi _{0}\right) =\frac{2\pi a\rho \left( \psi \right) 
}{v\left| \xi _{0}\right| }\sqrt{1+\left( \frac{R_{0}q\left( \psi \right) }{%
a\rho \left( \psi \right) }\right) ^{2}}\lambda \left( \psi ,\xi _{0}\right) 
\end{equation}
with 
\begin{equation}
\lambda \left( \psi ,\xi _{0}\right) =\left[ 1+\left( \frac{R_{0}q\left(
\psi \right) }{a\rho \left( \psi \right) }\right) ^{2}\right]
^{-1/2}\int_{\theta _{\min }}^{\theta _{\max }}\frac{d\theta }{2\pi }\frac{J%
}{a\rho \left( \psi \right) }\frac{\xi _{0}}{\xi }\frac{B}{\left|
B_{P}\right| }
\end{equation}

This complicated-looking expression is defined such as to take a very simple
intuitive form for the case of circular concentric flux-surfaces. See
Section ?? for details. It ensures that the normalized transit time of a
particle with only parallel velocity is exactly one.

\subsection{Idealized Magnetic Field Structure}

In a toroidal axisymmetric geometry, the magnetic field can be expressed
generally as 
\begin{equation}
\mathbf{B}=I\left( \psi \right) \nabla \phi +\nabla \psi \times \nabla \phi 
\end{equation}
so that 
\begin{eqnarray*}
B_{T} &=&I\left( \psi \right) \left| \nabla \phi \right|  \\
\left| B_{P}\right|  &=&\left| \nabla \psi \right| \left| \nabla \phi
\right| 
\end{eqnarray*}
and
\begin{eqnarray*}
B &=&\sqrt{B_{T}^{2}+B_{P}^{2}} \\
&=&\left| \nabla \phi \right| \sqrt{I^{2}\left( \psi \right) +\left| \nabla
\psi \right| ^{2}}
\end{eqnarray*}

We then have the ratio
\[
\frac{B}{\left| B_{P}\right| }=\sqrt{1+\frac{I^{2}\left( \psi \right) }{%
\left| \nabla \psi \right| ^{2}}}
\]
and we obtain
\[
\lambda \left( \psi \right) =\sqrt{1+\frac{I^{2}\left( \psi \right) }{\left|
\nabla \psi \right| ^{2}}}\left[ 1+\left( \frac{R_{0}q\left( \psi \right) }{%
a\rho \left( \psi \right) }\right) ^{2}\right] ^{-1/2}\int_{\theta _{\min
}}^{\theta _{\max }}\frac{d\theta }{2\pi }\frac{J}{a\rho \left( \psi \right) 
}\frac{\xi _{0}}{\xi }
\]

\subsection{Passing - Trapped particles}

The turning points are 
\begin{equation}
\theta _{\min }=\left| 
\begin{array}{cc}
-\pi & \text{for passing particles} \\ 
\theta _{T\min } & \text{for trapped particles}
\end{array}
\right.
\end{equation}
\begin{equation}
\theta _{\max }=\left| 
\begin{array}{cc}
\pi & \text{for passing particles} \\ 
\theta _{T\max } & \text{for trapped particles}
\end{array}
\right.
\end{equation}
where the trapping condition is given by $\left| \xi _{0}\right| <\xi _{0T}$
, where $\xi _{0T}$ is the pitch angle, defined at the minimum $B_{\min
}\left( \psi \right) $ on a given flux-surface, such that the parallel
velocity of the particle vanishes at the maximum $B_{\max }\left( \psi
\right) $. This can be expressed as a function of the two adiabatic
invariants 
\begin{eqnarray}
\mu &=&\frac{p_{\bot }^{2}}{2mB} \\
p^{2} &=&p_{\bot }^{2}+p_{\Vert }^{2}
\end{eqnarray}
such that 
\begin{eqnarray}
p_{\bot 1}^{2}+p_{\Vert 1}^{2} &=&p_{\bot 0}^{2}+p_{\Vert 0}^{2} \\
\frac{p_{\bot 1}^{2}}{B_{\max }} &=&\frac{p_{\bot 0}^{2}}{B_{\min }}
\end{eqnarray}

Then, setting $p_{\Vert 1T}=0$, we get 
\begin{equation}
p_{\Vert 0T}^{2}=\left( \frac{B_{\max }}{B_{\min }}-1\right) p_{\bot 0T}^{2}
\end{equation}
\begin{equation}
\xi _{0T}^{2}=1-\frac{B_{\min }}{B_{\max }}
\end{equation}

We can determine the turning angles $\theta _{T\min }$ and $\theta _{T\max }$
from the same system, so that, for a given $\left| \xi _{0}\right| <\xi
_{0T} $, 
\begin{equation}
\xi _{0}^{2}=1-\frac{B_{\min }}{B_{T\max }}
\end{equation}
so that 
\begin{eqnarray}
\theta _{T\min } &=&\theta \left( B=B_{T\max }\left| \theta <0\right. \right)
\\
\theta _{T\max } &=&\theta \left( B=B_{T\max }\left| \theta >0\right. \right)
\end{eqnarray}

\subsection{Case of Circular, Concentric Flux-Surfaces}

In this case, we have 
\begin{equation}
R=R_{0}+r\cos \theta 
\end{equation}
so that 
\[
\rho =\frac{r}{a}
\]
and therefore the toroidal field is 
\begin{equation}
B_{T}=B_{T0}\left( r\right) \frac{R_{0}}{R}=\frac{B_{T0}\left( r\right) }{
1+r/R_{0}\cos \theta }
\end{equation}
where 
\[
B_{T0}\left( r\right) =\frac{I\left( r\right) }{R_{0}}
\]
the poloidal field is 
\begin{equation}
B_{P}=B_{P0}\left( r\right) \frac{R_{0}}{R}=\frac{B_{P0}\left( r\right) }{%
1+r/R_{0}\cos \theta }
\end{equation}
where
\[
B_{P0}\left( r\right) =\frac{\left| \nabla \psi \right| }{R_{0}}
\]
and the total field is 
\begin{equation}
B=\frac{\sqrt{B_{P0}^{2}\left( r\right) +B_{T0}^{2}}}{1+r/R_{0}\cos \theta }=%
\frac{B_{0}}{1+r/R_{0}\cos \theta }
\end{equation}
with 
\begin{equation}
B_{0}=\sqrt{B_{P0}^{2}\left( r\right) +B_{T0}^{2}}
\end{equation}

Then, we have the ratio
\[
\frac{B}{B_{P}}=\frac{B_{0}}{B_{P0}\left( r\right) }=\sqrt{1+\frac{B_{T0}^{2}%
}{B_{P0}^{2}\left( r\right) }}
\]

We also have 
\begin{equation}
J=\frac{ds}{d\theta }=r=a\rho 
\end{equation}
and 
\begin{equation}
q\left( r\right) =\frac{r}{R_{0}}\frac{B_{T}}{B_{P}}=\frac{r}{R_{0}}\frac{%
B_{T0}}{B_{P0}}
\end{equation}
so that 
\[
\frac{B}{B_{P}}=\sqrt{1+\left( \frac{qR_{0}}{r}\right) ^{2}}
\]

The normalized bounce time then becomes simply
\begin{equation}
\lambda \left( \psi \right) =\int_{\theta _{\min }}^{\theta _{\max }}\frac{%
d\theta }{2\pi }\frac{\xi _{0}}{\xi }
\end{equation}

In addition, defining 
\begin{equation}
\epsilon =\frac{r}{R_{0}}
\end{equation}
we find 
\begin{eqnarray}
B_{\min } &=&\frac{B_{0}}{1+\epsilon } \\
B_{\max } &=&\frac{B_{0}}{1-\epsilon }
\end{eqnarray}
so that 
\begin{equation}
\xi _{0T}^{2}=\frac{2\epsilon }{1+\epsilon }
\end{equation}

The turning angles are determined by 
\begin{equation}
\xi _{0}^{2}=\frac{\left( 1-\cos \theta _{T}\right) \epsilon }{1+\epsilon }
\end{equation}
so that 
\begin{equation}
\theta _{T}=\arccos \left[ 1-\frac{\xi _{0}^{2}}{2\xi _{0T}^{2}}\right]
\end{equation}
\begin{eqnarray}
\theta _{T\min } &=&-\theta _{T} \\
\theta _{T\max } &=&\theta _{T}
\end{eqnarray}

Of course, the position of minimum $B$-field is 
\begin{equation}
\theta _{0}=0
\end{equation}

\end{document}
